\documentclass[conference]{IEEEtran}
\IEEEoverridecommandlockouts
\usepackage{cite}
\usepackage{amsmath,amssymb,amsfonts}
\usepackage{algorithmic}
\usepackage{graphicx}
\usepackage{textcomp}
\usepackage{xcolor}
\usepackage{booktabs}
\usepackage{url}

\def\BibTeX{{\rm B\kern-.05em{\sc i\kern-.025em b}\kern-.08em
    T\kern-.1667em\lower.7ex\hbox{E}\kern-.125emX}}

\begin{document}

\title{Water Quality Prediction via Gradient Boosting and Deep Neural Networks: Phase 1 Machine Learning Evaluation\\
\thanks{Applied Machine Learning (AML) Course Project -- Phase 1 Evaluation Deliverable.}
}

\author{\IEEEauthorblockN{AML Research Team}
\IEEEauthorblockA{\textit{Department of Computer Science \& Engineering} \\
\textit{Applied Machine Learning Laboratory} \\
Water Quality Telemetry Group}
}

\maketitle

\begin{abstract}
Automated water quality prediction from physical and chemical multi-sensor telemetry is essential for real-time environmental protection and public health compliance. In this work (Phase 1), we implement a mathematically rigorous Applied Machine Learning (AML) benchmark on the national Canadian Council of Ministers of the Environment (CCME) water monitoring repository comprising over 2.82 million observations. We prioritize physics-informed sensor data cleaning by flagging impossible thermodynamic values (e.g., pH outside 0--14 and thermal sensor spikes) using non-destructive fault indicators rather than row dropping. We conduct a thorough literature review identifying the foundational inductive bias gap between decision tree ensembles and deep neural networks: tree models partition input spaces via axis-aligned orthogonal cuts, which creates high-variance staircases when modeling smooth continuous chemical reaction manifolds. Following exploratory data analysis across 12 waterbody types, we construct domain chemical features (including temperature-dependent dissolved oxygen solubility polynomials and nitrogen speciation fractions) while enforcing a strict zero-leakage protocol and split-isolated imputation. Baseline models (Random Forest and XGBoost) are evaluated with 5-fold cross-validation across both continuous index regression and 5-class discrete quality tiers. XGBoost demonstrates superior continuous WQI prediction ($R^2 = 0.9992$, $\text{RMSE} = 0.4926$), while Random Forest excels in discrete classification ($\text{Accuracy} = 99.29\%$, $\text{Macro F1} = 0.9857$). Model calibration is verified via Expected Calibration Error (ECE), and coalitional game-theoretic TreeSHAP is used to uncover chemical drivers. This establishes an empirical baseline for Phase 2 deep neural network integration.
\end{abstract}

\begin{IEEEkeywords}
Data Cleaning, Water Quality Index, Gradient Boosting, XGBoost, Random Forest, Expected Calibration Error, Data Leakage, TreeSHAP, Inductive Bias.
\end{IEEEkeywords}

\section{Introduction}
Access to potable water is a critical imperative under the United Nations Sustainable Development Goal 6 (Clean Water and Sanitation). Traditional water quality determination relies on laboratory bioassays (such as 5-day Biochemical Oxygen Demand ($BOD_5$)), which introduce multi-day latency that prevents timely intervention during contamination surges. Modern environmental telemetry deploys continuous in-situ physical and chemical sensors (pH, dissolved oxygen, temperature, conductivity, nitrogen, and phosphorus compounds).

However, raw telemetry streams contain sensor calibration drift, hardware anomalies, and missing values. Phase 1 of this project establishes a disciplined, end-to-end Machine Learning evaluation framework that addresses data cleaning, literature grounding, exploratory data analysis, domain feature engineering, theoretical model formulations, cross-validated empirical benchmarking, and game-theoretic explainability.

\section{Data Cleaning \& Sensor Physical Bounds Auditing}
\subsection{Thermodynamic & Physical Bounds}
Before analyzing distributions or training estimators, telemetry must be audited against fundamental thermodynamic and chemical bounds:
\begin{enumerate}
    \item \textbf{pH Limits}: Aqueous solutions under standard environmental conditions are bounded by $\text{pH} \in [0.0, 14.0]$. Raw telemetry contained anomalous spikes up to $\text{pH} = 28.8$, indicating hardware electrode drift.
    \item \textbf{Water Temperature}: Surface water in natural temperate waterbodies exists in liquid phase within $[-5.0^\circ\text{C}, 50.0^\circ\text{C}]$. Uncalibrated thermal readings reached $97.5^\circ\text{C}$.
    \item \textbf{Dissolved Oxygen (DO)}: Dissolved oxygen concentration cannot be negative ($\text{DO} \ge 0.0\text{ mg/L}$). Negative values represent optical sensor zero-offset drift.
    \item \textbf{Detection Limits}: Concentrations of Ammonia, Orthophosphate, Nitrate, and BOD reported as $0.0\text{ mg/L}$ represent left-censored detection limits.
\end{enumerate}

\subsection{Non-Destructive Cleaning Protocol}
Dropping entire rows upon detecting an anomalous sensor reading discards valid co-occurring chemical measurements, reducing statistical power and biasing regional representations. Conversely, arbitrary value clipping conceals sensor failure from downstream estimators. 

We implement a non-destructive protocol (`flag\_and\_clean\_sensor\_anomalies`): out-of-bounds readings are replaced with $\text{NaN}$ for split-aware imputation, and an explicit binary flag (`[sensor]\_is\_sensor\_fault = 1`) is appended to the feature matrix. This signals hardware failure directly to tree splits.

\section{Literature Review}
\subsection{Tree-Based Ensembles on Tabular Data}
Breiman \cite{breiman2001random} established that aggregating decorrelated decision trees trained on bootstrap samples reduces model variance without increasing bias. Chen and Guestrin \cite{chen2016xgboost} introduced XGBoost, formalizing gradient tree boosting with second-order Taylor expansions of the objective function, exact greedy split finding, and explicit leaf-weight regularization ($\ell_1$ and $\ell_2$). Lundberg et al. \cite{lundberg2020local} introduced TreeSHAP, providing exact local and global feature attributions based on cooperative game theory that eliminate the cardinality bias of Gini impurity.

\subsection{The "Hybrid" Logic: Tabular Trees to Deep Neural Networks}
A foundational question in tabular machine learning is why tree ensembles traditionally outperform neural networks on tables, and where the boundary lies:
\begin{itemize}
    \item \textbf{Tree Model Inductive Bias}: Decision trees partition feature space using axis-aligned orthogonal cuts ($x_j \le \theta$). This enables rapid fitting of unscaled tabular data and step-function thresholds.
    \item \textbf{The Chemical Manifold Gap}: Chemical reaction equilibria, such as oxygen solubility curves and multi-nutrient buffering, lie along smooth, continuous diagonal manifolds. Approximating these smooth manifolds with orthogonal partitions forces tree models to construct deep, high-variance staircases of splits.
    \item \textbf{Phase 2 Motivation}: A regularized Multi-Layer Perceptron (MLP) with continuous nonlinear activations naturally parameterizes smooth multi-sensor chemical interaction manifolds, motivating the hybrid comparison in Phase 2.
\end{itemize}

\section{Dataset Quality \& Exploratory Data Analysis (EDA)}
We analyze the national Canadian Surface Water Quality Monitoring dataset, comprising 2,827,977 observation records across 61,689 geographic monitoring areas (`Area`) and 12 distinct waterbody classes (River: 65.5\%, Effluent: 21.3\%, Lake: 5.4\%, Estuarine: 1.7\%).

The dataset provides dual target definitions:
\begin{enumerate}
    \item \textbf{Continuous Target (`CCME\_Values`)}: Continuous Water Quality Index score ranging from $31.3$ to $100.0$ (Mean: 85.05, Median: 90.60).
    \item \textbf{Discrete Target (`CCME\_WQI`)}: Categorical quality tiers defined by the Canadian Council of Ministers of the Environment: Excellent (36.1\%), Good (36.5\%), Marginal (11.6\%), Fair (11.2\%), and Poor (4.5\%).
\end{enumerate}

Bivariate correlation analysis demonstrates that no individual sensor displays a correlation $|r| > 0.75$ with the target (Orthophosphate: $r = -0.74$, Nitrogen: $r = -0.61$, Ammonia: $r = -0.39$, BOD: $r = -0.32$), confirming that water quality represents multi-parameter chemical kinetics requiring multivariate nonlinear modeling.

\section{Feature Engineering \& Zero-Leakage Pipeline}
\subsection{Domain Chemical Interactions}
To aid tree models in capturing continuous chemical kinetics, we engineer domain features:
\begin{enumerate}
    \item \textbf{DO Saturation Ratio \& Deficit}: Theoretical oxygen solubility at 1 atm as a function of temperature ($T$ in $^\circ\text{C}$):
    \begin{equation}
    \text{DO}_{\text{sat}}(T) \approx 14.652 - 0.41022 T + 0.00799 T^2 - 0.000078 T^3
    \end{equation}
    The saturation percentage $\text{Ratio}_{\text{DO}} = \text{DO} / \text{DO}_{\text{sat}}(T)$ and oxygen deficit $\text{Deficit}_{\text{DO}} = \text{DO}_{\text{sat}}(T) - \text{DO}$ serve as indicators of eutrophication and organic decomposition.
    \item \textbf{Organic Stress Ratio}: Ratio of Biochemical Oxygen Demand to Dissolved Oxygen ($\text{BOD} / (\text{DO} + 10^{-3})$).
    \item \textbf{Nitrogen Speciation Fractions}: Proportions of toxic un-ionized ammonia and oxidized nitrate relative to total nitrogen ($\text{Ammonia} / (\text{Nitrogen} + 10^{-3})$).
    \item \textbf{Nutrient Enrichment Product}: $\text{Nitrate} \times \text{Orthophosphate}$ (eutrophication accelerator).
\end{enumerate}

\subsection{Data Leakage Prevention}
Deriving a target index from components and including those same components without auditing causes trivial algebraic memorization. We institute an automated correlation audit: any feature exhibiting Pearson $|r| \ge 0.98$ with the target is flagged. Zero data leakage was verified.

\subsection{Split-Aware Imputation}
Imputing missing values prior to train/test partitioning introduces lookahead bias. Our `SplitAwareImputer` fits median estimators strictly on $X_{\text{train}}$ and transforms $X_{\text{test}}$, appending missingness indicators to capture potential Missing-Not-At-Random (MNAR) patterns.

\section{Theoretical Rigor of Baseline Models}
\subsection{Random Forest (Variance Reduction)}
For an ensemble of $B$ trees with pairwise correlation $\rho$ and individual variance $\sigma^2$, the variance of the ensemble mean prediction $\bar{f}(x)$ is:
\begin{equation}
\text{Var}(\bar{f}(x)) = \rho \sigma^2 + \frac{1 - \rho}{B} \sigma^2
\end{equation}
Bootstrap sampling and random subspace feature selection ($m \approx \sqrt{p}$) minimize $\rho$, driving the second term toward zero as $B \to \infty$.

\subsection{XGBoost (Regularized Gradient Boosting)}
At iteration $t$, XGBoost minimizes the second-order Taylor expansion of the objective:
\begin{equation}
\tilde{\mathcal{L}}^{(t)} \approx \sum_{i=1}^n \left[ g_i f_t(x_i) + \frac{1}{2} h_i f_t^2(x_i) \right] + \gamma T + \frac{1}{2} \lambda \sum_{j=1}^T w_j^2
\end{equation}
where $g_i = \partial_{\hat{y}^{(t-1)}} l(y_i, \hat{y}^{(t-1)})$ and $h_i = \partial^2_{\hat{y}^{(t-1)}} l(y_i, \hat{y}^{(t-1)})$. The optimal weight for leaf $j$ containing instance set $I_j$ is given analytically by:
\begin{equation}
w_j^* = -\frac{\sum_{i \in I_j} g_i}{\sum_{i \in I_j} h_i + \lambda}
\end{equation}
Regularization parameter $\lambda$ dampens leaf weights, while $\gamma$ penalizes tree complexity, guarding against tabular overfitting.

\section{Model Application, Calibration \& Results}
The benchmark was trained on 50,000 stratified samples with an 80/20 train/test split and 5-fold cross-validation. Hyperparameters were optimized via randomized search.

\begin{table}[htbp]
\caption{Continuous WQI Regression Benchmark (Holdout Test Set)}
\begin{center}
\begin{tabular}{lcccc}
\toprule
\textbf{Model} & \textbf{MAE} & \textbf{RMSE} & \textbf{$R^2$} & \textbf{MAPE} \\
\midrule
Dummy (Mean Baseline) & 13.921 & 17.446 & -0.0001 & 20.47\% \\
Random Forest (Breiman) & 0.2544 & 0.7306 & 0.9982 & 0.36\% \\
\textbf{XGBoost (Chen \& Guestrin)} & \textbf{0.1520} & \textbf{0.4926} & \textbf{0.9992} & \textbf{0.22\%} \\
\bottomrule
\end{tabular}
\label{tab:regression_results}
\end{center}
\end{table}

\begin{table}[htbp]
\caption{Discrete WQI Classification Benchmark (5 Quality Classes)}
\begin{center}
\begin{tabular}{lcccc}
\toprule
\textbf{Model} & \textbf{Accuracy} & \textbf{Bal. Acc.} & \textbf{Macro F1} & \textbf{Weighted F1} \\
\midrule
Dummy (Stratified) & 29.09\% & 19.72\% & 0.1972 & 0.2909 \\
\textbf{Random Forest} & \textbf{99.29\%} & \textbf{98.54\%} & \textbf{0.9857} & \textbf{0.9929} \\
XGBoost & 98.84\% & 97.87\% & 0.9813 & 0.9884 \\
\bottomrule
\end{tabular}
\label{tab:classification_results}
\end{center}
\end{table}

XGBoost achieved an $R^2 = 0.9992$ and $\text{RMSE} = 0.4926$ on continuous regression, while Random Forest achieved $99.29\%$ accuracy and $0.9857$ Macro F1 on discrete quality classification. Probabilistic calibration was assessed via Expected Calibration Error (ECE), confirming tight confidence alignment on reliability diagrams.

\section{Interpretability via TreeSHAP}
To explain model decisions without Gini impurity bias, we compute exact Shapley values using TreeSHAP \cite{lundberg2020local}:
\begin{equation}
\phi_i(x) = \sum_{S \subseteq F \setminus \{i\}} \frac{|S|!(|F| - |S| - 1)!}{|F|!} [f_x(S \cup \{i\}) - f_x(S)]
\end{equation}
TreeSHAP summary beeswarm and bar plots identify Orthophosphate, Nitrogen, and Ammonia as primary negative drivers of the water quality index, with abnormal pH excursions and dissolved oxygen deficits acting as secondary tipping points.

\section{GitHub Repository \& Code Quality Architecture}
The codebase is structured as a modular Python package:
\begin{itemize}
    \item `src/data/`: Data loading (`loader.py`) and non-destructive sensor cleaning (`clean.py`).
    \item `src/features/`: Domain chemical feature engineering and leakage audits (`engineering.py`).
    \item `src/models/`: Baseline training (`train_baseline.py`), metrics (`evaluate.py`), and TreeSHAP explainability (`explainability.py`).
    \item `notebooks/`: Four dedicated Jupyter notebooks (`01_data_cleaning`, `02_dataset_quality_and_eda`, `03_feature_engineering`, `04_model_application`).
    \item `tests/`: Automated unit test suite with 100\% pass rate under `pytest`.
\end{itemize}

\section{Conclusion \& Phase 2 Roadmap}
Phase 1 establishes a comprehensive, mathematically grounded tabular machine learning benchmark. We have eliminated data leakage, instituted sensor physical validation, trained regularized Random Forest and XGBoost architectures, and quantified both discriminative performance and probabilistic calibration. In Phase 2, we will leverage this baseline to evaluate a small regularized Multi-Layer Perceptron (MLP) with continuous activations, directly testing the hypothesis that deep representations can improve upon tree ensembles when modeling continuous multi-sensor chemical interactions.

\begin{thebibliography}{00}
\bibitem{breiman2001random} L. Breiman, ``Random forests,'' \textit{Machine Learning}, vol. 45, no. 1, pp. 5--32, 2001.
\bibitem{chen2016xgboost} T. Chen and C. Guestrin, ``XGBoost: A scalable tree boosting system,'' in \textit{Proc. ACM SIGKDD Int. Conf. Knowl. Discov. Data Min.}, 2016, pp. 785--794.
\bibitem{lundberg2020local} S. M. Lundberg et al., ``From local explanations to global understanding with explainable AI for trees,'' \textit{Nature Machine Intelligence}, vol. 2, no. 1, pp. 56--67, 2020.
\bibitem{ccme2017} Canadian Council of Ministers of the Environment, ``CCME Water Quality Index 1.0 User's Manual,'' Canadian Environmental Quality Guidelines, 2017.
\bibitem{goodfellow2016deep} I. Goodfellow, Y. Bengio, and A. Courville, \textit{Deep Learning}, MIT Press, 2016.
\end{thebibliography}

\end{document}
